% Build: cd docs/reports/reconstruction/methods && latexmk -xelatex pdc_reco_literature_and_air_ms_gap_20260426.tex
\documentclass[a4paper, 11pt]{article}

\usepackage[margin=2.4cm]{geometry}
\usepackage{amsmath,amssymb}
\usepackage{booktabs}
\usepackage{multirow}
\usepackage{siunitx}
\usepackage{xcolor}
\usepackage{hyperref}
\usepackage{enumitem}
\usepackage{float}
\usepackage{graphicx}
\usepackage{tikz}
\usetikzlibrary{arrows.meta,positioning,calc}
\usepackage{longtable}
\usepackage{tabularx}

\definecolor{goodgreen}{RGB}{34,139,34}
\definecolor{warnred}{RGB}{200,50,50}

\hypersetup{
  colorlinks=true,
  linkcolor=blue!60!black,
  citecolor=blue!60!black,
  urlcolor=blue!60!black,
}

\title{\textbf{PDC Reconstruction Literature Survey \\ and SAMURAI Terminal Air-MS Gap Analysis}}
\author{Baiting Tian \\ RIKEN Nishina Center / DPOL Experiment}
\date{2026-04-26}

\begin{document}
\maketitle
\tableofcontents
\clearpage

\section{TL;DR}
This document surveys the published reference frame for the SAMURAI PDC1/PDC2 drift-chamber chain and contrasts it with our current Geant4 reconstruction performance, in order to localise the air multiple-scattering gap that drives our PDC2 spatial resolution one order of magnitude above design.

\begin{itemize}
  \item \textbf{Reference frame.} The PDC1/PDC2 chain has no dedicated peer-reviewed publication. The canonical references are the SAMURAI overview paper by Kobayashi et al., Nucl. Instr. Meth. B \textbf{317}, 294--304 (2013) §3.4\footnote{\label{fn:kobayashi2013}T.~Kobayashi et al., Nucl. Instr. Meth. B \textbf{317}, 294--304 (2013), DOI: \href{https://doi.org/10.1016/j.nimb.2013.05.089}{10.1016/j.nimb.2013.05.089}.}, the RIKEN internal technical note \texttt{Detector-PDC.pdf}\footnote{\label{fn:detectorpdc}RIKEN Nishina Center, \emph{Detector-PDC} technical note, \url{https://www.nishina.riken.jp/ribf/SAMURAI/image/Detector-PDC.pdf}.}, and the Kahlbow PhD thesis (TU Darmstadt, 2019)\footnote{\label{fn:kahlbow2019}J.~Kahlbow, PhD thesis, TU Darmstadt (2019), TUprints 9192, \url{https://tuprints.ulb.tu-darmstadt.de/9192/}.} which documents an independent reconstruction of the same hardware.
  \item \textbf{Design specifications.} Combining Kobayashi 2013 §3.4\footref{fn:kobayashi2013} with the SAMURAI 2012 construction proposal\footnote{\label{fn:samurai2012}SAMURAI Collaboration, \emph{SAMURAI Construction Proposal} (2012), \url{https://ribf.riken.jp/SAMURAI/120425_SAMURAIConstProp.pdf}.}, the PDC1/PDC2 chain targets 1~mm position resolution, 2~mrad angular resolution, and 0.8\% relative momentum resolution ($\Delta p/p$). Commissioning RI-beam runs reached $\Delta p/p \approx 1/1500$\footref{fn:kobayashi2013}; the original proton-mode design value is $\Delta p/p \approx 1/700$ as quoted in the RIBF-TAC05 (2005) construction proposal\footnote{\label{fn:ribftac05}SAMURAI Collaboration, \emph{RIBF-TAC05 SAMURAI construction proposal} (2005), \url{https://ribf.riken.jp/RIBF-TAC05/10_SAMURAI.pdf}.}.
  \item \textbf{Our measured $\sigma_y$(PDC2).} In our Geant4 + reconstruction chain we observe $\sigma_y(\text{PDC2}) = 12.83$--$14.93$~mm in the all-air configuration, $6.80$--$8.21$~mm when Region~A (dipole cavity + beam pipe) is set to vacuum (\texttt{beamline\_vacuum}), and approximately $1.0$~mm in the all-vacuum configuration\footnote{\label{fn:msablation}See \texttt{docs/reports/reconstruction/momentum\_reco/ms\_ablation/ms\_ablation\_mixed\_20260422\_zh.md} §4.}. The realistic-air values sit roughly an order of magnitude away from the 1~mm hardware design specification\footref{fn:kobayashi2013}.
  \item \textbf{Air MS dominates and the configuration choice matters.} A quadrature decomposition\footref{fn:msablation} attributes the all-air budget to $\sigma_A \approx 11.19$~mm from Region~A (cavity + beam pipe, $\sim$4.3~m) and $\sigma_B \approx 7.43$~mm from Region~B (Exit Window $\to$ PDCs, $\sim$2~m). The stage~D inputs from the same memo\footnote{See \texttt{docs/reports/reconstruction/momentum\_reco/ms\_ablation/ms\_ablation\_mixed\_20260422\_zh.md} §6.} give $\sigma_y(\text{PDC2}) = 7.49$~mm if Region~A is vacuum and $\sigma_y(\text{PDC2}) = 13.50$~mm if both Region~A and Region~B are air. Whether Region~A is air or vacuum during the experiment is therefore a configuration question that selects between these two values; this document does not assert which one corresponds to the actual run.
\end{itemize}

See §6 for proposed action items.

\section{The problem: air multiple scattering at the SAMURAI terminal}
% [content from Task 3]
\subsection{Geometry of the proton flight path}
\subsection{Measured \texorpdfstring{$\sigma$}{sigma} values}
\subsection{Gap to design specifications}
\subsection{Why this matters for the science floor}

\section{PDC reconstruction reference frame}
% [content from Tasks 4--6]
\subsection{Hardware foundations}
\subsection{Algorithm spine}
\subsection{Performance achieved across SAMURAI campaigns}

\section{Gap analysis}
% [content from Tasks 7--8]
\subsection{Performance: \texorpdfstring{$\sigma$}{sigma} values}
\subsection{Methodology}
\subsection{Material budget and vacuum configuration}
\subsection{Experimental configuration and geometry}

\section{Broader MINOS+SEASTAR chain context}
% [content from Task 9]
\subsection{SEASTAR programme timeline}
\subsection{Companion detectors}
\subsection{Reviews and resources}
\subsection{Caveats from the broader bibliography}

\section{Implications and next steps}
% [content from Task 10]
\subsection{Short-term data acquisition action items}
\subsection{Medium-term simulation and analysis directions}
\subsection{Open questions}
\subsection{Non-goals}

\section{References}
% [content from Task 11]
\subsection{Hardware and methodology (peer-reviewed, foundational)}
\subsection{Algorithm references}
\subsection{PhD theses with PDC methodology}
\subsection{Physics papers using PDC1/PDC2 (confirmed/implicit)}
\subsection{Excluded / context-only}
\subsection{Reviews and topical collections}
\subsection{Construction proposals and internal notes}
\subsection{Project-internal references}

\end{document}
